%% ============================================================
%% nomi_personaggi.tex — Registro dei nomi di PG e NPC
%% Vespri 1282
%%
%% Tutti i nomi sono storicamente attestati in Sicilia
%% nel XIII secolo, tratti da:
%% — Caracausi, Dizionario onomastico della Sicilia (1993)
%% — Marrone, Repertori del Regno di Sicilia 1282–1377
%% — Documenti della R. Cancelleria aragonese post-1282
%% — Cronache: Bartolomeo di Neocastro, Saba Malaspina
%% ============================================================

%% ============================================================
%% FAZIONE I — CONGIURATI
%% ============================================================

%% 1. Il Messaggero
%%    Ruggero di Amato — nome latino-normanno
%%    Ruggero: frequentissimo in Sicilia normanna (Ruggero I, II)
%%    Amato: cognome-patronimico attestato nei registri siciliani

%% 2. Il Cospiratore
%%    Gualtiero de Caltagirone — nome normanno + origine geografica
%%    Gualtiero (Gautier): normanno, attestato in molti atti del XIII sec.
%%    Caltagirone: toponimo siciliano-arabo (Qal'at al-Jiran)

%% 3. La Spia
%%    Margherita di Brindisi — nome greco-latino femminile
%%    Margherita: attestato nel regno svevo (Margherita di Navarra, regina)
%%    di Brindisi: patronimico geografico comune per agenti del sud

%% 4. Il Mercante
%%    Bartolomeo Lancia — nome latino + cognome svevo-siciliano
%%    Lancia: famiglia nobile attestata a corte di Federico II

%% 5. Il Soldato Disertore
%%    Drogo de Metz — nome franco + provenienza continentale
%%    Drogo: nome franco-normanno attestato nei registri militari angioini
%%    de Metz: indica origine lotaringia, plausibile per soldato di Carlo

%% 6. Il Nobile in Esilio
%%    Tancredi di Chiaramonte — nome normanno + casata storica
%%    Tancredi: re di Sicilia (1189–1194), nome simbolicamente carico
%%    Chiaramonte: famiglia baronale siciliana attestata dal XII sec.

%% 7. Il Falsario
%%    Niccolò il Cretese — nome greco + epiteto d'origine
%%    Niccolò: greco, diffusissimo in Sicilia (influenza bizantina)
%%    il Cretese: Palermo ospitava una comunità greca di Creta

%% ============================================================
%% FAZIONE II — BARONI
%% ============================================================

%% 8. Il Barone
%%    Simone di Mazzarino — nome greco-latino + feudo storico
%%    Simone: attestato frequentemente nei registri baronali siciliani
%%    Mazzarino: feudo reale nel Nisseno, attestato nel XIII sec.

%% 9. Il Cavaliere
%%    Riccardo Bonello — nome normanno + cognome attestato
%%    Riccardo: normanno, attestato nei registri militari
%%    Bonello: cognome siciliano medievale (de Bonello, atti notarili)

%% 10. La Castellana
%%    Elisabetta di Paternò — nome latino + feudo etneo
%%    Elisabetta: attestato nell'aristocrazia femminile del XIII sec.
%%    Paternò: feudo etneo, attestato come possesso baronale

%% 11. Il Notaio
%%    Matteo di Termini — nome ebraico-latino + città portuale
%%    Matteo: diffusissimo nel clero e nella burocrazia medievale
%%    Termini: Termini Imerese, sede di notariato attestato

%% 12. Il Figlio
%%    Goffredo di Mazzarino — nome normanno-germanico
%%    Goffredo (Godefridus): attestato nei registri feudali siciliani
%%    [figlio di Simone — stessa casata]

%% 13. Il Medico (Baroni) — nota: ha una passione per la medicina occulta,
%%    contatti cross-fazione: Yusuf al-Balarmì (Clero) e Maimone (Popolo)
%%    Maestro Filippo Carafa da Salerno — titolo + nome + provenienza
%%    Filippo: greco, attestato tra i medici del regno
%%    da Salerno: la Scuola Medica Salernitana era il riferimento del XIII sec.

%% 14. La Moglie del Traditore
%%    Costanza Ventimiglia — nome svevo + casata siciliana
%%    Costanza: nome simbolico (Costanza d'Altavilla, madre di Federico II)
%%    Ventimiglia: famiglia baronale siciliana attestata dal XII sec.

%% ============================================================
%% FAZIONE III — CLERO
%% ============================================================

%% 15. Il Vescovo
%%    Berardo di Castronovo — nome germanico + feudo ecclesiastico
%%    Berardo: arcivescovo di Palermo (1206–1244) — nome dal forte peso
%%    Castronovo: diocesi siciliana, feudo ecclesiastico attestato

%% 16. Il Frate Predicatore
%%    Fra' Tommaso di Caltagirone — nome latino + origine geografica
%%    Tommaso: dopo Tommaso d'Aquino (†1274), nome diffusissimo nei Domenicani
%%    Caltagirone: città con forte presenza domenicana nel XIII sec.

%% 17. Il Canonico
%%    Pandolfo di Messina — nome longobardo + città
%%    Pandolfo: attestato nei capitoli cattedrale siciliani del XIII sec.
%%    di Messina: canonici messinesi frequenti nella cattedrale di Palermo

%% 18. La Badessa
%%    Suor Adalasia di Troina — nome normanno + città storica
%%    Adalasia (Adelaide): moglie di Ruggero I, nome normanno di rilievo
%%    Troina: prima diocesi normanna di Sicilia

%% 19. Il Sacrista
%%    Arnaldo Gennarino — nome catalano-normanno + cognome siciliano
%%    Arnaldo: nome d'origine catalana, plausibile nel contesto pre-aragonese
%%    Gennarino: attestato nei registri delle chiese palermitane

%% 20. Il Medico Arabo
%%    Yusuf ibn Khalaf al-Balarmì — nome arabo + nisba palermitana
%%    Yusuf (Giuseppe): nome coranico diffusissimo
%%    ibn Khalaf: patronimico ('figlio di Khalaf')
%%    al-Balarmì: nisba = 'il Palermitano' (Balarm = Palermo in arabo)

%% 21. Il Pellegrino
%%    Luca di Bari — nome greco-evangelico + città pugliese
%%    Luca: attestato tra i pellegrini nei registri delle ospitalità medievali
%%    di Bari: Bari era punto di partenza per pellegrinaggi verso la Terra Santa

%% ============================================================
%% FAZIONE IV — POPOLO
%% ============================================================

%% 22. L'Artigiano
%%    Calogero Ferrante — nome greco-siciliano + cognome di mestiere
%%    Calogero: dal greco kalos-geron ('bel vecchio'), tipicamente siciliano
%%    Ferrante: da ferrus/faber ferrarius, attestato tra gli artigiani

%% 23. La Pescivendola
%%    Perna di Cassaro — nome latino + quartiere
%%    Perna: attestato nei registri femminili siciliani del XIII sec.
%%    di Cassaro: indicazione del quartiere d'appartenenza

%% 24. Il Ragazzo
%%    Salvo Russo — nome latino + cognome di colore/aspetto
%%    Salvo: dal latino salvus, attestato nel repertorio siciliano medievale
%%    Russo: da russus ('rosso'), cognome frequente in Sicilia

%% 25. Il Pescatore
%%    Pietro di Cefalu — nome apostolico + città costiera
%%    Pietro: diffusissimo, spesso tra i pescatori (richiamo evangelico)
%%    di Cefalù: città portuale siciliana, pescatori attestati nei registri

%% 26. La Vedova
%%    Giacoma la Nera — nome latino + soprannome
%%    Giacoma: femminile di Giacomo, attestato nella Sicilia del XIII sec.
%%    la Nera: soprannome per capelli/carnagione scura, frequente

%% 27. Il Guaritore di Strada
%%    Maimone di Palermo — nome arabo sicilianizzato + città
%%    Maimone: arabismo (Maymun = 'fortunato'), frequente nella Sicilia
%%    cristiana post-islamica anche tra non-musulmani

%% 28. Il Vecchio
%%    Basilio di Monreale — nome greco + città normanna
%%    Basilio: greco, attestato tra i cristiani di rito greco-siculo
%%    di Monreale: la cattedrale normanna ospitava comunità di varia origine

%% ============================================================
%% NPC TRASVERSALI
%% ============================================================

%% Jean de Saint-Rémy — nome già attestato storicamente (manteniamo)
%% Arcivescovo di Palermo — nel 1282: Giovanni di Laon (manteniamo il ruolo,
%%   il GM può usare il nome storico o un nome fittizio)
%%   Per il gioco: Guglielmo di Monforte (plausibile per un presule del XIII sec.)

%% ============================================================
%% NPC DEI CONGIURATI
%% ============================================================

%% Emissario di Pietro d'Aragona
%%    Guillem de Montcada — nome catalano-aragonese
%%    Guillem (Guglielmo): forma catalana, attestato tra i nobili aragonesi
%%    Montcada: famiglia nobile catalana presente nei registri siciliani post-1282

%% Contatto di Messina
%%    Enrico di Taormina — nome germanico + città ionica
%%    Enrico: diffuso nel XIII sec. (Enrico VI, padre di Federico II)
%%    di Taormina: città con attiva rete mercantile

%% Agente Bizantino
%%    Alexios Philanthropenos — nome greco-bizantino
%%    Alexios: frequente nella burocrazia imperiale di Costantinopoli
%%    Philanthropenos: famiglia aristocratica bizantina attestata nel XIII sec.

%% ============================================================
%% NPC DEI BARONI
%% ============================================================

%% Vassallo Filo-Angioino
%%    Ugo di Santapau — nome normanno + feudo
%%    Ugo (Hugo): normanno, attestato tra i feudatari filo-angioini
%%    Santapau: feudo siciliano (Licodia Eubea) attestato nel XIII sec.

%% Capofamiglia Rivale
%%    Ruggero di Incisa — nome normanno + feudo minore
%%    Ruggero: normanno classico
%%    di Incisa: feudo minore nel palermitano, attestato nei registri

%% ============================================================
%% NPC DEL CLERO
%% ============================================================

%% Fra' Matteo — nome già assegnato nel manuale (manteniamo)
%%    Fra' Matteo di Sant'Agata — aggiunta del patronimico geografico
%%    Sant'Agata: paese siciliano, plausibile origine per un frate minore

%% Legato Pontificio
%%    Gherardo da Parma — nome italiano settentrionale + provenienza
%%    Gherardo: attestato tra i legati pontifici del XIII sec.
%%    da Parma: i legati erano spesso italiani settentrionali fidati a Roma

%% Monaca Visionaria (Suor Eufemia) — nome già nel manuale, manteniamo
%%    Suor Eufemia di Corleone — aggiunta patronimico
%%    Corleone: città interna siciliana, sede di conventi femminili

%% ============================================================
%% NPC DEL POPOLO
%% ============================================================

%% Soldato Francese
%%    Renaud de Coucy — nome franco + casata piccola-nobiliare
%%    Renaud: normanno-franco, comune tra i soldati di Carlo
%%    de Coucy: famiglia piccolo-nobiliare dell'Ile-de-France

%% Capoquartiere
%%    Salamone dell'Albergheria — nome ebraico-siciliano + quartiere
%%    Salamone: attestato tra i cristiani siciliani (dal re Salomone)
%%    dell'Albergheria: quartiere popolare di Palermo

%% La Madre del Morto
%%    Addolorata di Seralcadio — nome devozionale + quartiere
%%    Addolorata: attestato nei registri femminili siciliani
%%    di Seralcadio: quartiere artigiano di Palermo
